\section{Historischer Abriss und Originalarbeiten}

\textbf{Euler: Summieren und Differenzenbilden als Gegenoperationen.}
In den \emph{Institutiones calculi differentialis} von 1755 behandelt
Euler endliche Differenzen und ihre Anwendung auf Folgen. In Teil~I,
Kapitel~II, \S\S~53--54 bildet er zu einer Folge ihre Partialsummenfolge,
die \emph{series summatrix}. Die Differenzen benachbarter Partialsummen
geben die ursprünglichen Glieder zurück. Euler setzt ausdrücklich eine
Null vor die Partialsummenfolge, um auch das erste Glied zu erfassen
\cite{abel:euler}. Das ist die Rechenidee hinter dem Zusammenspiel
von Summation und Differenzenbildung. Eine erstmalige Entdeckung des
Teleskopierens wird damit nicht behauptet.

\textbf{Abel: Die Umformung dient der Abschätzung.}
Abels Untersuchung der Binomialreihe erschien 1826 im ersten Band
von Crelles Journal. Abschnitt~II, Lehrsatz~III führt die entscheidende
Rechnung aus: Abel schreibt die Glieder $t_k$ als Differenzen ihrer
Partialsummen $p_k$ und gruppiert die mit $\varepsilon_k$ gewichtete
Summe neu. Es entstehen die Differenzen
$\varepsilon_k-\varepsilon_{k+1}$ und der Randterm
$p_m\varepsilon_m$ \cite{abel:original}. In unserer Schreibweise
ist das (A).

Abel verwendet die Identität unmittelbar für eine einseitige
Abschätzung bei positiven abnehmenden Gewichten. Der folgende
Lehrsatz~IV wendet sie auf Potenzreihen an. Der ursprüngliche Nutzen
entspricht damit dem Kern dieser Lesefassung: Eine Kontrolle der
Partialsumme wird auf gewichtete Summen übertragen. Unsere
Betragsabschätzung (E) ist die dazu passende zweiseitige Form.

Der Originalartikel von 1826 ist deutsch. Die gut zugängliche
französische Fassung in der Werkausgabe von 1839 enthält die
Umformung unter \emph{Théorème~III} auf Druckseite~69 und die
anschließende Anwendung unter \emph{Théorème~IV} auf S.~69--70.
Diese Ausgabe ist ein Wiederabdruck in französischer Sprache und
nicht der Originaldruck von 1826.

\textbf{Dirichlet: Konvergenz durch Auslöschung.}
Dirichlets Arbeit über trigonometrische Reihen von 1829 untersucht
oszillierende Integrale. Im Beweis zerlegt er ein solches Integral
in Abschnitte mit wechselnden Vorzeichen und abnehmenden Beträgen
\cite{abel:dirichlet}. Das zeigt den analytischen Anwendungshorizont
der Auslöschung. Die oben gegebene moderne Herleitung des
Dirichlet-Kriteriums für Reihen ist keine Wiedergabe dieses
Fourier-Beweises. Auch der Beweis von (A) aus der Produktidentität
ist hier auf einen möglichst kurzen heutigen Beweisgang zugeschnitten.

\begin{thebibliography}{9}
\small
\setlength{\itemsep}{.65em}

\bibitem{abel:euler}
Leonhard Euler:
\emph{Institutiones calculi differentialis cum eius usu in
analysi finitorum ac doctrina serierum}.
Sankt Petersburg / Berlin, 1755.
Pars~I, Caput~II, \S\S~53--54.
\href{https://doi.org/10.3931/e-rara-3788}{Originaldigitalisat der ETH-Bibliothek}.
Zum leichteren Lesen:
\href{https://17centurymaths.com/contents/euler/diffcal/part1ch2.pdf}{Ian Bruces Ausgabe mit lateinischem Text und englischer Übersetzung, Teil~I, Kapitel~II}.
Die Paragraphenangaben beziehen sich auf Eulers Text.

\bibitem{abel:original}
Niels Henrik Abel:
\emph{Untersuchungen über die Reihe
$1+\frac m1x+\frac{m(m-1)}{1\cdot2}x^2+
\frac{m(m-1)(m-2)}{1\cdot2\cdot3}x^3+\cdots$ u.s.w.}
Journal für die reine und angewandte Mathematik~1 (1826),
S.~311--339; hier Abschnitt~II, Lehrsatz~III, S.~314,
und der anschließende Lehrsatz~IV.
\href{https://doi.org/10.1515/crll.1826.1.311}{Dauerhafter Nachweis des Originalartikels};
\href{https://gdz.sub.uni-goettingen.de/id/PPN243919689_0001}{Originalband im Göttinger Digitalisierungszentrum}.
Französische Fassung: \emph{Recherche sur la série \ldots},
in B.~Holmboe (Hg.), \emph{Œuvres complètes de N.~H.~Abel},
Bd.~I, Christiania, 1839, S.~66--92;
hier Abschnitt~II, Théorème~III, S.~69, und Théorème~IV, S.~69--70.
\href{https://www.hist-math.fr/textes/Abel1826_TheoremeBinomial.pdf}{Digitalisat dieser Werkausgabe}.

\bibitem{abel:dirichlet}
Peter Gustav Lejeune Dirichlet:
\emph{Sur la convergence des séries trigonométriques qui servent
à représenter une fonction arbitraire entre des limites données}.
Journal für die reine und angewandte Mathematik~4 (1829), S.~157--169.
\href{https://geodesic.mathdoc.fr/item/JRAM_1829__4_183134/}{Bibliografischer Nachweis bei Mathdoc};
\href{https://arxiv.org/abs/0806.1294}{Transkription des Originaltextes durch Ralf Stephan (2008)}.
Die Transkription berichtigt einzelne Druckfehler des Originals.

\end{thebibliography}
